\documentclass[12pt]{article}

% Idioma e codificação
\usepackage[brazil]{babel}
\usepackage[utf8]{inputenc}
\usepackage[T1]{fontenc}

% Matemática e listas
\usepackage{amsmath, amssymb}
\usepackage{enumitem}

% Layout e cores
\usepackage{geometry}
\usepackage{graphicx}
\usepackage{xcolor}

% Cores institucionais
\definecolor{maincolor}{RGB}{0,100,50}
\definecolor{forestgreen}{RGB}{0,100,50}

% Cabeçalho e rodapé
\usepackage{fancyhdr}

% Blocos
\usepackage[most]{tcolorbox}
\newtcolorbox{block}[1]{%
  colback=green!5,
  colframe=maincolor,
  coltitle=white,
  title=#1,
  fonttitle=\bfseries,
  boxrule=1.5pt,
  arc=3mm,
  left=6pt,
  right=6pt,
  top=6pt,
  bottom=6pt
}

% Margens
\geometry{margin=2.5cm}

% Fancyhdr
\pagestyle{fancy}
\fancyhf{}

\renewcommand{\headrulewidth}{2pt}
\renewcommand{\footrulewidth}{2pt}

\renewcommand{\headrule}{\hbox to\headwidth{%
  \color{forestgreen}\leaders\hrule height \headrulewidth\hfill}}

\renewcommand{\footrule}{\hbox to\headwidth{%
  \color{forestgreen}\leaders\hrule height \footrulewidth\hfill}}

% Altura do cabeçalho
\setlength{\headheight}{52pt}
\addtolength{\topmargin}{-20pt}

% Cabeçalho
\fancyhead[L]{%
  \textbf{UniRV} \\ \small Universidade de Rio Verde
}

\fancyhead[R]{%
  \raggedleft
  \textcolor{forestgreen}{%
    \small
    \textbf{Lista de Exercícios: Monitoração SRE}\\
    Me.~Sergio Souza Novak\\
    \textit{Engenharia da Qualidade e Confiabilidade}
    }
}
    
% Rodapé
\fancyfoot[C]{\textcolor{forestgreen}{\thepage}}

% Comando para caixas de resposta
\newtcolorbox{resposta}[1]{%
  colback=white,
  colframe=forestgreen!30,
  arc=1mm,
  height=#1,
  width=\textwidth,
  boxrule=0.8pt,
  breakable
}

% =========================
\begin{document}
% =========================

\begin{center}
  {\Large \textbf{Lista de Exercícios: Monitoração Agregada e Princípios SRE}}\\[0.2cm]
  \textit{Engenharia de Software / Engenharia da Qualidade}
\end{center}

\vspace{0.4cm}

\section*{Exercícios de Fixação}

\textbf{1.} De acordo com os princípios do Google SRE, defina monitoração e explique por que ela é considerada a base da Hierarquia de Necessidades de Confiabilidade.
\begin{resposta}{3.5cm}
\end{resposta}

\textbf{2.} Os sistemas de monitoração produzem diferentes saídas. Explique a diferença conceitual e de urgência entre \textbf{Alertas (Paging)}, \textbf{Tickets} e \textbf{Logging}.
\begin{resposta}{4.5cm}
\end{resposta}

\textbf{3.} Diferencie a monitoração de \textbf{Caixa-Branca (White-box)} da monitoração de \textbf{Caixa-Preta (Black-box)}. Em qual situação cada uma é mais indicada para um SRE?
\begin{resposta}{4cm}
\end{resposta}

\textbf{4.} O Google define os "4 Sinais de Ouro" da monitoração. Liste-os e forneça uma breve definição para cada um, exemplificando sua aplicação em um sistema HTTP.
\begin{resposta}{5.5cm}
\end{resposta}

\textbf{5.} Sobre a métrica de \textbf{Latência}, por que é fundamental monitorar separadamente a latência de requisições que resultaram em erro das que resultaram em sucesso?
\begin{resposta}{3.5cm}
\end{resposta}

\textbf{6.} Explique os conceitos de \textbf{SLI}, \textbf{SLO} e \textbf{SLA}. Como eles se relacionam hierarquicamente na gestão de serviços?
\begin{resposta}{4cm}
\end{resposta}

\textbf{7.} Se um serviço possui um SLO de disponibilidade de $99,95\%$ em um período mensal, qual é o seu \textbf{Orçamento de Erro (Error Budget)} em termos percentuais? O que deve acontecer se esse orçamento for totalmente consumido antes do fim do período?
\begin{resposta}{4cm}
\end{resposta}

\textbf{8.} No contexto de alertas, qual a diferença entre um alerta baseado em \textbf{Causa} e um alerta baseado em \textbf{Sintoma}? Por que o Google recomenda priorizar alertas baseados em sintomas com impacto no usuário?
\begin{resposta}{4.5cm}
\end{resposta}

\textbf{9.} Por que a \textbf{Média} pode ser enganosa ao analisar a latência de um sistema distribuído? Qual a vantagem de utilizar \textbf{Percentis} (como P95 ou P99) em vez da média aritmética?
\begin{resposta}{4cm}
\end{resposta}

\textbf{10.} Um sistema web registrou as seguintes latências de resposta (em milissegundos) durante um período de monitoração. A tabela apresenta cada valor observado e sua respectiva frequência.

\begin{center}
\begin{tabular}{c|c}
\textbf{Latência (ms)} & \textbf{Frequência} \\
\hline
100 & 15 \\
105 & 18 \\
110 & 22 \\
120 & 20 \\
150 & 12 \\
200 & 6 \\
800 & 3 \\
2500 & 2 \\
5000 & 1 \\
12000 & 1 \\
\end{tabular}
\end{center}

\begin{enumerate}[label=\alph*)]

\item Calcule a \textbf{média aritmética} das latências considerando suas frequências.

\item Determine a \textbf{mediana (P50)} da distribuição.

\item Compare os valores obtidos e responda:

\begin{itemize}
\item Qual medida representa melhor a experiência da \textbf{maioria dos usuários}?
\item Qual medida revela a existência de \textbf{latências extremas (cauda da distribuição)}?
\end{itemize}

\item Explique por que, em sistemas distribuídos e em engenharia de confiabilidade (SRE), percentis como \textbf{P95, P98 ou P99} costumam ser mais informativos do que a média.

\end{enumerate}

\begin{resposta}{15cm}
\end{resposta}
% =========================
\end{document}
% =========================
